\chapter{LLM 辅助文献阅读、结构化笔记与报告写作}

\section{本章导读：文献和报告中的 AI 输出必须回到来源}

LLM 很擅长帮助我们整理摘要、提取结构、比较观点和润色报告。但文献阅读与科研写作有一条不能让步的底线：最终 claim 必须回到可核查来源。AI 生成的摘要如果没有原文页码、图表、公式或上下文支撑，就不能直接变成报告结论。

本章把 LLM 用作阅读和写作流程中的结构化助手，而不是替代阅读的机器。它可以帮你建立笔记表、整理 claim ledger、提醒可能的反例、改写段落；但引用是否准确、文献是否被误读、结论是否超出证据，都需要人工核对。

读这一章时，请把每个重要句子都当成需要追踪来源的 claim。AI 可以加快整理速度，但不能让你跳过“我为什么能这样说”的证据责任。

\section{案例目标}

第 35 章讨论的是：当 AI 帮你写出一段科研代码时，怎样通过验证回路判断它值不值得留下。第 36 章再往前走一步，讨论了多步骤代理怎样在工具边界、停止条件和人工接管之内运行。

本章继续沿着这条主线前进，但把场景换成另一个本科生更快会遇到的问题：\textbf{当你开始用 LLM 帮你读论文、摘笔记、写报告时，怎样避免“读得很快，但记得不准、引得不清、写得过头”。}

本章的目标有三点：

\begin{enumerate}
    \item 理解为什么文献阅读场景里，LLM 更适合先帮助你建立结构化 claim ledger，而不是直接产出最终摘要；
    \item 学会把论文阅读拆成“证据卡片、claim 去重、citation / quote 核查、caveat 收集、人工复核”几个可审计步骤；
    \item 学会把阅读结果整理成一份更适合写进 notebook、周报或课程报告的最小结构化摘要。
\end{enumerate}

\section{为什么这一章紧跟在第 35 章和第 36 章之后}

如果说第 35 章强调的是“AI 生成代码之后要验证”，第 36 章强调的是“AI 运行多步骤流程时要有工具边界与 human handoff”，那么本章强调的就是另一个同样现实的问题：\textbf{AI 帮你读论文时，也不能跳过证据边界。}

很多学生第一次把 LLM 用到文献阅读时，会自然地提出类似请求：

\begin{itemize}
    \item “帮我总结这篇论文的主要贡献”；
    \item “帮我把结果写成组会汇报的三点结论”；
    \item “帮我写一个 related work 段落”。
\end{itemize}

这些请求本身没有错，但危险在于：如果你直接要最终 prose，而没有先要求它保留 claim、出处、caveat 和待核查项，那么它很容易把“讨论里的推测”“附录里的未核对数字”“还没确认页码的引用”一起包装成流畅结论。

所以本章的核心立场是：\textbf{LLM 可以显著加快文献阅读，但真正可靠的中间产物不是自由摘要，而是结构化、可回查的 claim ledger。}

\section{一个极小的论文证据卡片数据集}

配套 notebook 使用 \nolinkurl{literature_reading_workflow_demo.csv}。这不是一篇真实论文的全文，也不是在线抓取结果，而是一组完全本地、教学用途的\textbf{证据卡片}。每一行都代表学生在阅读论文时摘出的一条信息，包含：

\begin{itemize}
    \item \texttt{note\_id}；
    \item \texttt{section\_label}；
    \item \texttt{claim\_id}；
    \item \texttt{salience\_score}；
    \item \texttt{note\_role}；
    \item \texttt{citation\_key}；
    \item \texttt{citation\_status}；
    \item \texttt{reference\_route}；
    \item \texttt{human\_check\_required}；
    \item \texttt{excerpt\_text}。
\end{itemize}

这里的设计非常刻意。数据里既有可以直接保留的结果 claim，也有几类最容易让 AI 摘要出错的条目：

\begin{itemize}
    \item 两条高显著度、但其实属于同一个 claim 的重复结果；
    \item 一条讨论段里的外推性推测；
    \item 两条必须写进报告的 caveat；
    \item 一条需要做页码 / quote 级核对的数字；
    \item 一条最好保留在方法附注里的背景信息。
\end{itemize}

因此，本章真正要训练的不是“怎样让 AI 写得更像摘要”，而是更基础也更重要的能力：

\begin{itemize}
    \item 你能不能让 AI 帮你先整理证据，而不是先替你下结论？
    \item 你能不能区分“可直接写进结果”“必须连 caveat 一起写”“必须人工复核”“最好先不要写进报告”这几种不同状态？
\end{itemize}

\section{先看一个很常见的 baseline：直接按显著度抽三条}

如果没有结构化工作流，一个很常见的 baseline 就是：\textbf{按 salience score 取前三条，然后让 AI 直接把它们写成摘要。}

\begin{examplebox}[直接抽三条的 baseline]
\begin{lstlisting}[style=pythonstyle]
top_notes = sorted(rows, key=lambda row: row["salience_score"], reverse=True)[:3]
\end{lstlisting}
\end{examplebox}

这个做法的问题在文献阅读里非常常见：

\begin{itemize}
    \item 它可能重复选到同一个 claim 的两个表述，导致摘要看起来“信息很多”，但其实 claim coverage 很低；
    \item 它可能把讨论段里的推测当成结果；
    \item 它几乎一定会漏掉限制条件与 caveat；
    \item 它不会自动提醒你哪些数字还需要 quote / page check。
\end{itemize}

在本章教学数据上，这个 baseline 取到的 top-3 里会混入一条不该直接写进报告的推测性表述，而且还会重复同一个核心结果 claim。因此 notebook 中会看到表~\ref{tab:ch37_baseline_vs_ledger} 这样的最小对照：

\begin{table}[htbp]
\centering
\small
\setlength{\tabcolsep}{4pt}
\begin{tabular}{@{}p{0.22\textwidth}>{\centering\arraybackslash}p{0.17\textwidth}>{\centering\arraybackslash}p{0.17\textwidth}>{\centering\arraybackslash}p{0.16\textwidth}p{0.20\textwidth}@{}}
\toprule
方法 & 结果 claim precision & claim coverage & caveat coverage & 教学解读 \\
\midrule
按显著度直接抽三条 & 0.67 & 0.50 & 0.00 & 容易重复同一 claim，并把推测混进摘要 \\
claim ledger + route workflow & 1.00 & 1.00 & 1.00 & 先去重、分流，再决定哪些句子能进入报告 \\
\bottomrule
\end{tabular}
\caption{本章 notebook 中两个最小文献阅读流程的对照。重点不是“谁写得更像摘要”，而是谁更尊重证据边界。}
\label{tab:ch37_baseline_vs_ledger}
\end{table}

\section{claim ledger 的关键不是“摘要更长”，而是“证据状态更清楚”}

本章 notebook 把阅读过程拆成几个非常具体的步骤：

\begin{itemize}
    \item \textbf{claim 去重}：把属于同一个 \texttt{claim\_id} 的多条摘录先归并；
    \item \textbf{route 分流}：把每条证据卡片送进“可直接写结果”“必须连 caveat 保留”“需要人工复核”“方法附注”或“暂不写入报告”这几类桶；
    \item \textbf{citation / quote check}：任何 \texttt{citation\_status != ok} 的条目都不直接写进最终 prose；
    \item \textbf{caveat 收集}：限制条件不能因为“摘要篇幅有限”就被自动丢掉；
    \item \textbf{report skeleton}：最后才把通过核查的结果 claim 和 caveat 拼成报告骨架。
\end{itemize}

最小 routing 逻辑可以写成：

\begin{examplebox}[一个最小文献阅读 routing 逻辑]
\begin{lstlisting}[style=pythonstyle]
if reference_route == "exclude":
    exclude
elif citation_status != "ok" or human_check_required == "yes":
    human_review
elif reference_route == "caveat":
    caveat
elif reference_route == "verified_result":
    verified_result
else:
    method_note
\end{lstlisting}
\end{examplebox}

请注意，这里的重点不是让 AI“自由发挥阅读理解”，而是让它在你给定的证据结构里\textbf{先保留证据状态，再生成文字。}

\section{为什么 caveat 不是摘要的附属品}

学生在写 reading note 时最常见的错误之一，就是只保留“好消息”：

\begin{itemize}
    \item AUC 提升了多少；
    \item purity 提升了多少；
    \item 哪个方法在基准测试里更好。
\end{itemize}

但论文阅读真正可靠的地方，往往恰恰在 caveat 里：

\begin{itemize}
    \item 提升是否只发生在高 S/N 样本上；
    \item 训练数据是否只来自单一 survey；
    \item 讨论里的外推是否真的被结果支撑；
    \item 某个数字是不是还需要你自己回到原文核页码。
\end{itemize}

因此，本章 notebook 会强制把 caveat 作为独立输出，而不是让它们依赖“AI 有没有顺手提到”。这和第 35 章、第 36 章的逻辑是连续的：\textbf{让 AI 帮你更快整理内容，但边界条件一定要显式保留下来。}

\section{先有 claim ledger，再写报告 skeleton}

很多人一开始就要求 AI “帮我写成一段 report prose”。更可靠的顺序其实相反：

\begin{enumerate}
    \item 先建立证据卡片；
    \item 再做 claim 去重；
    \item 再把每条 note route 到不同桶里；
    \item 最后才让 AI 帮你把结果整理成段落。
\end{enumerate}

也就是说，更适合交给 AI 的不是“一步写完报告”，而是：

\begin{itemize}
    \item 帮你把 claim ledger 变整齐；
    \item 帮你比较两个 claim 是否重复；
    \item 帮你把 caveat 单独列出来；
    \item 帮你把 human review 队列写成一个更清楚的核查清单；
    \item 等这些都稳定之后，再帮你润色 prose。
\end{itemize}

\section{一个更适合文献阅读的 prompt 结构}

如果任务已经进入论文阅读和报告写作，prompt 不应该只有“帮我总结”。更好的做法是把\textbf{claim、证据出处、caveat 和待核查项}一起写清楚。例如：

\begin{examplebox}[一个更适合科研文献阅读的 prompt 模板]
\begin{lstlisting}[style=pythonstyle]
prompt = """
你是我的科研文献阅读助手。

任务：
1. 先阅读证据卡片表，而不是直接写最终摘要；
2. 输出一个 claim ledger，区分 verified result、caveat、
   human review、method note 和 exclude；
3. 对每个 claim 保留 note_id、citation 状态和是否需要
   quote/page-level check；
4. 只有在 claim ledger 完整后，才生成报告 skeleton。

输出格式：
- Verified claims
- Caveats
- Human review queue
- Method notes
- Draft report skeleton
"""
\end{lstlisting}
\end{examplebox}

这里真正起作用的，不是什么神秘的提示词技巧，而是你有没有先要求 AI 保留\textbf{证据结构}。如果这个结构一开始就没有建立好，后面的 prose 再流畅，也只是在放大不透明的错误。

\section{配套 notebook 做了什么}

本章 notebook 采用的是一个 teaching-first 的最小设计：

\begin{itemize}
    \item 先读取证据卡片表，并打印 section、note role、citation 状态统计；
    \item 用按显著度抽三条的方式做一个 freeform-summary baseline；
    \item 再实现 claim 去重与 routing workflow；
    \item 输出 verified claims、caveats、human review queue 和 method notes；
    \item 最后打印一个更适合文献阅读的 prompt 模板、workflow checklist 和 reading snapshot。
\end{itemize}

这份 notebook 不直接连接在线论文 API，也不追求真实 RAG 系统。它真正要教的是：\textbf{在读论文这件事上，AI 最应该帮助你整理证据，而不是跳过证据。}

\section{和前面章节、以及 Capstone 的关系}

到这里为止，Part V 其实已经形成了一条很清楚的现代 AI 工作流主线：

\begin{itemize}
    \item 第 30 章到第 34 章讲模型怎样表示函数、局部模式、低维表示与长程上下文；
    \item 第 35 章讲 AI 生成科研代码后怎样验证；
    \item 第 36 章讲多步骤代理流程怎样被约束和审计；
    \item 本章讲 AI 进入论文阅读与报告写作时，怎样保留证据与核查边界。
\end{itemize}

到了 Capstone 项目里，这条逻辑会非常实用。真实项目中，你很可能会让 AI 帮你：

\begin{itemize}
    \item 整理 paper notes；
    \item 对比多篇论文的结果与 caveat；
    \item 生成组会 summary；
    \item 起草项目报告中的 related work 和 discussion 初稿。
\end{itemize}

但这些任务能不能安全地自动化，不取决于 AI“写得像不像摘要”，而取决于你有没有保留 claim ledger、引用核查和 human review 这几层中间结构。

进入 Part VI 后，claim ledger 会支撑 Interpretation card 与 Integrity card：前者说明哪些结论有证据、哪些只是 caveat，后者说明引用、页码和人工复核是否已经完成。这样，AI 辅助写作就不会把证据链压扁成一段流畅但不可回查的文字。

\begin{warnbox}[AI 辅助文献阅读中的常见误区]
\begin{itemize}
    \item 一开始就要求 AI 直接写 final summary，而不先保留 claim ledger；
    \item 不区分结果、推测、方法说明和限制条件；
    \item 让未核查页码的数字直接进入报告；
    \item 把 caveat 当成“可有可无的补充”，而不是正文的一部分；
    \item 以为 prose 流畅就等于阅读准确。
\end{itemize}
\end{warnbox}

\section{AI 助手如何帮助你}

在这个案例里，AI 助手最适合帮助你：

\begin{itemize}
    \item 把分散的证据卡片整理成更清楚的 claim ledger；
    \item 标出重复 claim 和互相冲突的表述；
    \item 帮你把 caveat 单独整理成 review note；
    \item 帮你把 human review 队列改写成更适合手工核对的 checklist；
    \item 在结构稳定后，再帮你把报告骨架润色成更自然的 prose。
\end{itemize}

但最终哪些句子真的可以被你写进报告、哪些数字必须回到原文页码逐条核对、以及哪些讨论段外推并没有被结果充分支撑，仍然必须由你自己判断。

\section{练习}

\begin{exercisebox}[基础练习]
把教学数据里一条 \texttt{citation\_status == verify\_page} 的 note 改成 \texttt{ok}，重新运行 notebook，观察 human review queue 和 final report skeleton 怎样变化。
\end{exercisebox}

\begin{exercisebox}[进阶练习]
请为本章 notebook 增加一个“跨论文对照”版本：再手工加入两条来自第二篇假想论文的证据卡片，并让 AI 帮你区分“共同结论”和“只在单篇论文中出现的 claim”。
\end{exercisebox}

\begin{exercisebox}[开放练习]
结合你自己的研究方向，设计一个“AI 阅读论文前必须保留的最小 claim ledger 模板”。至少包含：claim 本身、note / page 来源、引用状态、是否为结果或 caveat、以及是否需要人工复核。
\end{exercisebox}

\section{本章小结}

本章最重要的结论不是“怎样让 AI 更快帮你写摘要”，而是：\textbf{在科研文献阅读里，LLM 的真正价值不在于替你跳过阅读，而在于帮助你把证据、claim、引用状态和 caveat 整理得更清楚；真正让阅读结果可信的，仍然是你保留下来的 claim ledger、citation check 和 human review 结构。}
